In the European day-ahead electricity market, bids for each quarter-hour of the
next day must be submitted before the auction closes at 12:00.
Market participants use price forecasts to prepare these bids, so a forecast is
only operational if all of its inputs are available at that time. Many studies use the
load, solar, and wind forecasts published by the transmission system operators as
price-model inputs. The operators publish
these forecasts to meet a legal reporting obligation and do not disclose how they
are produced. The solar and wind forecasts are, in addition, only published after
12:00, so they could not be used in real operation.

This thesis develops an open and operational pipeline that produces its own load,
solar, and wind forecasts from open data and uses them to forecast day-ahead
prices for the German-Luxembourg market. It asks how accurate these forecasts are
compared with those of the system operators, whether they improve the price
forecast, and how the accuracy of the price forecast changes with information sets that grow hourly
from 07:00 to 12:00. Over a February to July 2026 test period, the
proposed load forecast reduces the error of the system operators' forecast by
36.0\%, while the proposed solar and wind forecasts are less accurate but
available in time. Using these forecasts lowers the price forecast error by
11.2\%. Forecasts made later in the morning become gradually more accurate, but the
only statistically significant gain comes at 12:00, when the price of the
Austrian day-ahead auction, which closes earlier, is already available. The results show that open
load, solar, and wind forecasts improve operational price forecasts even where
they are less accurate than those of the system operators. Beyond these results, the thesis provides
a transparent alternative to the forecasts of the system operators, whose methods
are not disclosed.
The pipeline uses only open data and is publicly released, and its forecasts are
submitted every day to the open Energy-Arena benchmark, where their accuracy can be
verified independently beyond the test period. Instead of relying on black boxes, anyone can inspect, rerun, and reuse
these benchmarked forecasts for their own purposes.
